\documentclass[pdflatex,sn-vancouver-num]{sn-jnl}

\usepackage{graphicx}%
\usexpackage{multirow}%
\usepackage{amsmath,amssymb,amsfonts}%
\usepackage{amsthm}%
\usepackage{mathrsfs}%
\usepackage[title]{appendix}%
\usepackage{xcolor}%
\usepackage{textcomp}%
\usepackage{manyfoot}%
\usepackage{booktabs}%
\usepackage{algorithm}%
\usepackage{algorithmicx}%
\usepackage{algpseudocode}%
\usepackage{listings}%
%\usepackage{paper}
%
%\usepackage{xstring}
\usepackage{csquotes}
\usepackage{pgfplots}
\usepackage{csvsimple}
\usepackage{listings}
\usepackage{siunitx}
%\usepackage{multirow}
%\pgfplotsset{compat=1.5}
%\bibliographystyle{bibliography}
%\usepgfplotslibrary{external}
%\tikzexternalize
\usepackage{xr}
\usepackage{todonotes}
\externaldocument{supplements}

\usepackage{array}
\newcolumntype{L}[1]{>{\raggedright\let\newline\\\arraybackslash\hspace{0pt}}m{#1}}
\newcolumntype{C}[1]{>{\centering\let\newline\\\arraybackslash\hspace{0pt}}m{#1}}
\newcolumntype{R}[1]{>{\raggedleft\let\newline\\\arraybackslash\hspace{0pt}}m{#1}}

%
% Workaround for local / relative path problems in conjunction with reading csv file.
%
%\newcommand{\csvpath}{}

%
% Workaround for faster compile without scatter plots:
%
%\newif\ifscatter
%\scattertrue

\newcommand{\taxon}{taxon}
\newcommand{\ku}{KrakenUniq}
\newcommand{\ktwo}{Kraken 2 (S)}
\newcommand{\ktwohc}{Kraken 2 (C)}
\newcommand{\ganon}{Ganon (S)}
\newcommand{\ganonlowfpr}{Ganon (C)}
\newcommand{\ganonall}{Ganon (A)}
\newcommand{\genestrip}{Genestrip (C)}
\newcommand{\genestriphs}{Genestrip (S)}
\newcommand{\hiseq}[1]{HiSeq}
\newcommand{\miseq}[1]{MiSeq}
\newcommand{\fastqone}[1]{Extr. from \cite{Wu:2024aa}}

\begin{document}

% Enter title:
    \title{Genestrip-FT: Optimized $k$-mer databases through refined taxonomies}

    \author*[1]{\sur{Daniel} \fnm{Pfeifer}}\email{daniel.pfeifer@hs-heilbronn.de}

    \author*[1]{\sur{Helai} \fnm{Zahelzai}}\email{hzahelzai@stud.hs-heilbronn.de}
    
    \author*[1]{\sur{Anton} \fnm{Atanasov}}\email{atanasov@stud.hs-heilbronn.de}
    
    \affil*[1]{\orgdiv{IT Faculty}, \orgname{Heilbronn University}, \orgaddress{\street{Max-Planck-Str. 39}, \city{Heilbronn}, \postcode{74081}, \state{Baden-W\"urttemberg}, \country{Germany}}}

% Enter date:
    \date{August 2026}


    \abstract{\textbf{Background:} 

    \textbf{Results:}
    
    \textbf{Conclusions:}     }

    \keywords{Metagenomics, $k$-mer counting, read classification, small $k$-mer database, high precision, high recall, high performance}

    \maketitle

    \section{Background}\label{background}

    \subsection{Introduction}
\label{problem}
    
    \begin{figure}[h]%
        \centering

         \caption{Box-plot histogram of relative $k$-mer distributions across ranks macro-averaged on a per species-basis for ?? different bacterial species. The distribution is based on all respective bacterial genomes in Refseq where $k$-mers shared with other any other bacterial species from RefSeq have been pushed up the ranks accordingly due to Genestrip's LCA update.}\label{kmerongenusborrelia}
    \end{figure}
    
    \begin{figure}[h]%
        \centering

         \caption{Histogram of relative $k$-mer distributions across species for the genus Borrelia (138) with ?? subordinate species. The distribution is based on all ?? $k$-mers of all Borrelia genomes in Refseq where such $k$-mers are shared between any two Borrelia species and have been pushed up at just to the rank genus due to Genestrip's LCA update.}\label{kmerongenusborrelia}
    \end{figure}

    \begin{figure}[h]%
        \centering

         \caption{Box-plot histogram of relative $k$-mer distributions across species on the rank genus averaged across ?? bacterial genera with an average of ?? subordinate species. The distribution is based on all ?? $k$-mers of respective genomes in Refseq where the $k$-mers have been pushed to the rank genus due to Genestrip's LCA update.}\label{kmeronbacterialgenera}
    \end{figure}
    
    \begin{figure}[h]%
        \centering

         \caption{Box-plot histogram of relative $k$-mer distribution across species on the rank genus averaged across ?? protozoan genera with ?? subordinate species. The distribution is based on all ?? $k$-mers of respective genomes in Refseq where the $k$-mers have been pushed to the rank genus due to Genestrip's LCA update.}\label{kmerviaranks}
    \end{figure}
    
    \subsection{Approach}
    \label{approach}

  \begin{figure}[h]%
        \centering

         \caption{}\label{lcasketch}
    \end{figure}   

  \begin{figure}[h]%
        \centering

         \caption{(1) to (3): Schematic representation of the lowest common ancestor update in a simplified $k$-mer database. (4): Genestrip-FT's taxonomic refinement the $k$-mer database from (3).}\label{ftsketch}
    \end{figure}   
   
   
    \section{Implementation}\label{methods}

  \begin{figure}[h]%
        \centering

         \caption{Schematic representation of the lowest common ancestor update in a simplified $k$-mer database.}\label{lcasketch}
    \end{figure}   

  \begin{figure}[h]%
        \centering

         \caption{Intersection counts between Borrelia species for $k$-mers pushed to the genus Borrelia via the LCA updated according to Figure \ref{lcasketch}}\label{intersectioncounts}
    \end{figure}
    
  \begin{figure}[h]%
        \centering

         \caption{Dendrogram of Borrelia species obtained via agglomerative clustering on the basis of Jaccard-index over the intersection counts from Figure \ref{intersectioncounts}}\label{dendrogram}
    \end{figure}
    

  \begin{figure}[h]%
        \centering

         \caption{Comparison of original database contents and the resulting refined database contents for Borrelia according to the dendrogram from Figure \ref{dendrogram}. The database contents is depicted as part of taxonomy tree extract in both cases.}\label{dendrogram}
    \end{figure}
    
    
    \section{Results}\label{results}
    
    \subsection{Intrinsic quality measures}
    
    In this section we measure the quality of the taxonomic refinement irrespective of read analysis and read classification. For this purpose, we define precision and recall only with respect to a $k$-mer database $d$ that has a taxonomy tree associated where $k$-mers have taxons, i.e. nodes, assigned from that tree.
    
    Let $t$ be a \emph{species taxon}, i.e. a taxon on the rank \enquote{species} with genomic files directly associated. For example in RefSeq, genomic files are usually associated with species taxons. We assume that $d$'s taxonomy tree is compact in sense that any leaf node in the tree represents a species taxon with one or more genomic files that contain at least one $k$-mer stored in $d$. If genomic files are associated with taxons below a species taxon $t$ and we simply consider these files as belonging to $t$ itself.\footnote{Note, that this is a simplification made for the adequacy of the intrinsic quality measures from below -- it is not a restriction that matters outside of this section.}
    
    \textbf{True positive plus false positive counts $TP + FP$ per species taxon $t$ in database $d$:} The number of unique $k$-mers in $d$ associated with any taxon on the path from $t$ to the taxonomic root including $t$ itself.
    
    \textbf{True positive counts $TP$ per species taxon $t$ in database $d$:} The number of unique $k$-mers in $d$ that occur in genomic files associated with $t$ where in $d$, such $k$-mers are associated with some taxon on the path from $t$ to the taxonomic root including $t$ itself.
    
    \textbf{False negative counts $FN$ per species taxon $t$ in database $d$:} The number of unique $k$-mers in $d$ that occur in genomic files associated with $t$ where in $d$, such $k$-mers are \emph{not} associated with any taxon on the path from $t$ to taxonomic root including $t$ itself.
    
    Based on these counts, the quality of the taxonomic refinement per species taxon $t$ in $d$ can be assessed via \emph{node precision} (equal to $TP / (TP + FP)$) and \emph{node recall} (equal to $TP / (TP + FN)$).
    We report average precision and average recall by averaging over all species taxons or subsets thereof. In particular, we also report averages over all species taxons under the same genus.
            
    Note that for a taxonomy tree with branching factor $\geq 2$ (which is the standard case) a high average precision ensures the a $k$-mer $h$ in $d$ is located as closely as possible above or right at species taxons whose genomic files contain $h$. Otherwise, if $h$ was located higher up in the taxonomy, then $h$ would also be on the path to the root starting from a species taxon $s$ whose genomic files do not contain $h$. But this would increase the $FP$-count of $s$ and in turn lower the average precision.
    
    Likewise, a high average recall ensures that a $k$-mer $h$ in $d$ is assigned to a tree node such that \emph{as many species taxons as possible} which contain $h$ in their genomic files are subordinate or equal to that node. Otherwise, the $k$-mer $h$ would not be on the path to the root for such a taxon. But this would increase the $FN$-count of that taxon and in turn lower the average recall.
    
    So together, both, a high average precision and a high average recall suggest that $k$-mers are located at a near optimal position in the taxonomy tree -- not too high up to lower precision but also not too far down to lower recall.
    
    \subsection{Validation against other tools}
    
    \subsection{Improvement of classification quality}
    
    Refinement at each level and node of the taxonomy tree.
    
    Taxon-based positive counts versus path-based positive counts for read classification:
    
    Taxon-based positive count for rank $r$: $1$ if the class taxon is equal to the correct taxon for rank $r$. Else $0$.
    
    Path-based positive count for rank $r$: If the class taxon is on the path between the taxonomic root and the correct taxon for rank $r$: $1$ divided by the number of taxons of rank $r$ that are subordinate or identical to the class taxon for rank $r$. Else $0$.
        
    \subsection{Use for phylogenic analysis}
    
    \subsection{Principle arguments in favour of the approach}
    
    \subsubsection{Comparison with established phyloegeny trees}
    
    Example candidates: Borrelia, Babesia, Malaria, Orthoxosvirus
    
    \subsubsection{Database refinement rerformance}
    
    \section{Conclusions}\label{discussion}


\section{Availability and requirements}

\textbf{Project name:} Genestrip\\
\textbf{Project home page:} \url{https://github.com/pfeiferd/genestrip}\\
\textbf{Operating system(s):} Platform independent\\
\textbf{Programming language:} Java\\
\textbf{Execution requirements:} JRE 11 or higher \\
\textbf{Build requirements:} JDK 11 or higher and Maven 2 or 3\\
\textbf{License:} Apache 2.0 subject to the Commons Clause License Condition v1.0 \\
\textbf{Any restrictions to use by non-academics:} License needed for commercial use\\

\section{List of abbreviations}

\begin{tabular}{lp{10cm}}
$t$ & Placeholder for taxon \\
$f$ & Placeholder for fastq file \\
$d$ & Placeholder for $k$-mer database \\
$k_{t,f,d}$& Number of $k$-mers found for taxon $t$ in fastq file $f$ via database $d$\\
$u_{t,f,d}$& Number of unique $k$-mers found for taxon $t$ in fastq file $f$  via database $d$\\
$u_{t,d}$& Number of $k$-mers stored under taxon $t$ in database $d$ \\
% $r_{t,f,d}$ &$k$-mer consistency ratio for $k$-mers found for taxon $t$ in fastq file $f$ with respect to database $d$ \\
\texttt{cv\_ku} & Complete viral database for KrakenUniq \\
\texttt{cv\_g} & Complete viral database for Genestrip \\
\texttt{hv\_ku} & Human virus database for KrakenUniq \\
\texttt{hv\_g} & Human virus database for Genestrip \\
FPP & False positive probability of a probabilistic filter \\
BPs & Base pairs \\
\texttt{MicrobialDB} & Large, comprehensive microbial database for KrakenUniq \\
\texttt{tb} & Tick-borne bacterial pathogens database for Genestrip \\
DB & Database
\end{tabular}

\section{Declarations}

\subsection{Ethics approval and consent to participate}

Not applicable

\subsection{Consent for publication}

Not applicable

\subsection{Availability of data and materials}

The exact code base for this publication is available under\\
\url{https://github.com/pfeiferd/genestrip/releases/tag/v2.7}.\\

\subsection{Competing interests}

The authors declare that they have no competing interests.

\subsection{Funding}

The work behind this publication was not funded.

\subsection{Authors' contributions}

DP developed the software, designed the evaluation and wrote the manuscript. \todo{What else?}

\subsection{Acknowledgements}

\clearpage
    
\end{document}